\section{Related Work}
Early PIC studies established that spatial grid resolution and shape-function choice control aliasing noise and numerical heating~\cite{langdon1970,okuda1972}.
Brackbill and Barnes showed that discrete charge non-conservation can destabilize long-time integrations~\cite{brackbill1980}, motivating exact deposition schemes.
Esirkepov's trajectory-based method provides rigorous charge conservation for arbitrary shape factors~\cite{esirkepov2001}, while Eastwood's virtual-particle formulation extended conservative deposition to electromagnetic PIC~\cite{eastwood1981}.
Harvey et al.\ demonstrated charge-conserving deposition on hybrid CPU--GPU systems~\cite{harvey2015}.

Recent work has focused on adapting PIC algorithms to emerging architectures.
Viereck et al.\ surveyed memory-access and parallelization challenges for deposition on modern hardware~\cite{viereck2016}.
Wilke et al.\ reported GPU-accelerated PIC strategies for large-scale supercomputers~\cite{wilke2017}.
Our contribution differs by providing a unified, reproducible 2D reference implementation that separates sort cost from deposit-kernel throughput and couples CPU microbenchmarks with charge, energy, spectral-noise, and Langmuir-wave validation under a common runtime interface.
